% =====================================================================
%  sec5.tex — 5. Wild Ramification
%  原文 pp.7–10
% =====================================================================
\section{野分歧}

基于与数论的类比，几何 Langlands 通常并不只对驯顺分歧的情形表述。人们会进一步追问一个类似的、涉及带有任意阶极点处的全纯联络的丛 $E\to C$ 的对偶陈述。换言之，在把 $E$ 在一个点 $p\in C$ 附近的全纯结构平凡化之后，联络看起来像
\begin{equation}
A=dz\left(\frac{T_{n}}{z^{n}}+\frac{T_{n-1}}{z^{n-1}}+\cdots+\frac{T_{1}}{z}+\cdots\right),
\label{eq:wild}
\end{equation}
其中正则项被省略。极点阶数大于 $1$ 的亚纯联络，称为有非正则奇点。

要为此情形表述一个对偶陈述，乍看之下对此处所述的几何 Langlands 进路构成严峻挑战。我们的基本观点是：那些根本的对偶陈述只依赖 $C$ 作为一个定向二维流形。$C$ 上的复结构只是作为一种工具被引入，用来回答某些在不引入复结构的情况下也能提出的自然问题。

从这个观点看，驯顺分歧是自然的，因为这种意义上的单极点有清楚的拓扑含义。在点 $p\in C$ 处有单极点的亚纯联络，是编码 $C'=C\setminus p$ 上平坦联络绕 $p$ 的单值性的一种自然方式。而这一单值性当然是纯粹拓扑的概念。然而，一个阶数大于 $1$ 的极点，其拓扑含义可能是什么呢？

一个密切相关的观察是：$T_{1}$ 是 $A$ 在 $z=0$ 处极点的留数\zhnt{5.1}，因而与局部坐标 $z$ 的选择无关。然而高阶极点的系数 $T_{2},\dots,T_{n}$ 却肯定依赖于局部坐标的选择\zhnt{5.2}。我们怎能希望把它们纳入一个只应依赖 $C$ 作为定向二维流形的理论呢？

而且，如果计划是先表述一个拓扑性质的对偶猜想、再用 Hitchin 方程去证明它，那么我们就要面对 Hitchin 方程是否与非正则奇点相容的问题。关于对 $\Phi=(A,\phi)$ 的 Hitchin 方程，其形式概略地说是 $d\Phi+\Phi^{2}=0$。若在 $z=0$ 附近奇性满足 $|\Phi|\sim1/|z|^{n}$，则 $|d\Phi|\sim1/|z|^{n+1}$，而 $|\Phi|^{2}\sim1/|z|^{2n}$。当 $n=1$ 时，$d\Phi$ 与 $|\Phi|^{2}$ 量级相当，因此 Hitchin 方程看起来是合理的。然而当 $n>1$ 时，$|\Phi|^{2}\gg|d\Phi|$，Hitchin 方程中的非线性项看来会太强\zhnt{5.3}。

不过，两个问题都有自然的答案。对第一个问题的答案是：尽管表面上不然，人们确实可以给非正则奇点的联络关联上某种超出普通单值性、且具有纯拓扑含义的东西。恰当的概念是扩张单值性，它既包含 Stokes 矩阵，也包含普通单值性。Stokes 矩阵是常微分方程非正则奇点经典理论的一部分，例如见 Wasow (1965)\zhnt{5.4}。

为简洁起见，假定联络奇异部分的首项系数 $T_{n}$ 是正则半单的，那么可以作一个规范变换，把 $T_{1},\dots,T_{n}$ 共轭到极大环面中。然后定义模空间 $\YC(G,C;T_{1},\dots,T_{n})$，它在规范变换下参数化这样的对：$C$ 上的全纯 $G_{\C}$-丛，以及一个具有式 \eqref{eq:wild} 所述非正则奇点的联络。如 Boalch (2001) 所证明，这个空间 $\YC(G,C;T_{1},\dots,T_{n})$ 自然地是一个复辛流形，其复辛结构只依赖 $C$ 作为定向二维流形。这一点可以通过把 Atiyah 与 Bott (1982) 中表述的辛结构的规范场论定义搬到当前情形来证明。

此外，$\YC(G,C;T_{1},\dots,T_{n})$ 的复辛结构与 $T_{2},\dots,T_{n}$ 无关（只要 $T_{n}$ 保持半单）。

此时，Jimbo 等 (1981) 引入的重要概念——等单值形变——就登场了。存在一种自然的方式改变参数 $T_{2},\dots,T_{n}$，而不改变由 $\YC(G,C;T_{1},\dots,T_{n})$ 所参数化的广义单值数据。而且，如 Boalch (2001) 新近证明的，广义单值数据空间的复辛结构在等单值形变下不变\zhnt{5.5}。于是粗略地说，可以定义一个只依赖 $T_{1}$ 与整数 $n\ge1$ 的复辛流形 $\YC(G,C;T_{1},n)$。

参数 $T_{2},\dots,T_{n}$ 按刚才所述的意义是非本质的，这当然是受欢迎的，因为正如我们已经注意到的，它们没有明显的拓扑含义。现在我们处在与非分歧、驯顺分歧情形非常相似的局面中。给定复辛流形 $\YC(G,C;T_{1},n)$ 及其复辛形式 $\Omega$，我们可以定义它的 B-模型，或者用实辛形式 $\omega=\operatorname{Im}\Omega$ 定义它的 A-模型。当然，对偶群也可以同样处理，定义另一个复辛流形 $\YC(\LG,C;{}^{L}T_{1},n)$，及其自己的 B-模型与 A-模型。而正如非分歧情形一样，我们可以问这两个模型是否互为镜像。

甚至在尝试回答这个问题之前，我们应先把问题稍作精细化。由于要求 $T_{n}$ 是正则半单的，简单地把 $T_{n}$ 忘掉并不完全正确。当 $T_{n}$ 变化时可能有单值性。我们把 $T_{n}$ 看作取值于 $\mathfrak{t}_{\C}^{\mathrm{reg}}\otimes K_{p}^{n-1}$，记号如下：$\mathfrak{t}_{\C}$ 是 $G_{\C}$ 中极大环面的李代数，$\mathfrak{t}_{\C}^{\mathrm{reg}}$ 是其中由正则元构成的子空间，$K_{p}$ 是 $C$ 的余切丛在 $p$ 处的纤维。$\mathfrak{t}_{\C}^{\mathrm{reg}}$ 的基本群称为 $G$ 的辫群；我们记之为 $B(G)$\zhnt{5.6}。由于这些单值性，为了定义 $\YC(\LG,C;{}^{L}T_{1},n)$，实际上需要选取基点 $*\in\mathfrak{t}_{\C}^{\mathrm{reg}}$；更精确地说，我们可以把这个空间记作 $\YC(\LG,C;{}^{L}T_{1},n,*)$。群 $B(G)$ 通过单值性同时作用在 $\YC(G,C;T_{1},n,*)$ 的 B-模型与 A-模型上。对偶地，相应的辫群 $B(\LG)$ 作用在 $\YC(\LG,C;{}^{L}T_{1},n,*)$ 的 B-模型与 A-模型上。

然而，$B(G)$ 与 $B(\LG)$ 这两个群自然同构；确实，模去一个不变二次型的选择，存在从 $\mathfrak{t}_{\C}^{\mathrm{reg}}$ 到 ${}^{L}\mathfrak{t}_{\C}^{\mathrm{reg}}$ 的自然映射，因此这两个空间有相同的基本群，且其中一个中基点的选择（在同伦意义下）决定了另一个中的基点。一个更好（但仍不精确）的问题是：$\YC(G,C;T_{1},n,*)$ 与 $\YC(\LG,C;{}^{L}T_{1},n,*)$ 之间是否存在一个与辫群作用交换的镜像对称。

我们还期待这一镜像对称只依赖 $C$ 作为定向二维流形，从而与映射类群交换。我们可以把 $C$ 的映射类群与辫群看作扮演相当平行的角色。事实上，由于上一段中典范丛的纤维 $K_{p}$ 的出现，这两个群并不只是彼此交换；作用着的群是映射类群被 $B(G)$ 的一个扩张\zhnt{5.7}。

正如驯顺分歧情形一样，要得到正确的镜像对称猜想，我们需要把参数稍作扩充，以得到一个镜像对称的参数组。但我们也面对 Hitchin 方程是否与野分歧相容这一根本问题。如上所述，Hitchin 方程的非线性使这一点乍看可疑。但幸运的是，结果表明一切正常，如 Biquard 与 Boalch (2004) 所证明。要点是：同样在假定 $T_{n}$ 正则半单的前提下，我们可以假定联络的奇异部分是交换的。Hitchin 方程虽然非线性，在交换情形下却变成线性的；而一旦交换化，它们就与任意阶的奇性相容。以此出发，结果表明对任何 $n$ 都可以发展出一套与非分歧情形中更熟悉的故事相当平行的、带非正则奇点的 Hitchin 方程理论。例如，方程解的模空间 $\MH$ 是超凯勒的。在其中一个复结构下，$\MH$ 参数化具有类似式 \eqref{eq:wild} 那种奇性的平坦联络；在另一个复结构下，它参数化 Higgs 丛 $(E,\phi)$，其中 $\phi$ 有阶数为 $n$ 的类似极点。存在 Hitchin 纤维化，所有通常的性质都有密切的类比。

这一切恰好给出了正确的素材，可用 Hitchin 方程来建立这两个模空间之间所期望的镜像对称。详细的解释（其中这一经典几何被嵌入四维规范场论）见 Witten (2007)。许多论证与前面给出的论证相当类似。
